\begin{center}
    {\Large\itshape Alban au début du XVII\ieme{} siècle}
\end{center}

\vspace{1em}

Le cadastre de 1598 (ou compoix) fait état des biens possédés
par N.D. de Pietat. Comme cela vient en tête du registre,
avant les biens des hospitaliers, et ceux des Passemar de
Saint André, il est possible de penser que c'est là le vocable
de N.D. d'Alban. Ces appellations de N.D. de la Merci ou de
la Pitié sont fort à la mode depuis le XV\ieme{} siècle.
Ainsi la cathédrale d'Albi possède en ces temps-là une
chapelle de N.D. de Pitié... dans la collégiale de Saint
Sernin également. Ce peut n'être aussi que le nom d'une
confrérie. Quant aux hospitaliers de Saint Jean de Jérusalem,
héritiers des templiers, ces religieux, comme on sait, ont
essaimé un peu partout chez nous. Leur chef, dans le coin,
est le Commandeur de Rayssac. Mais il y a des hospitaliers
à Carménel (l'escalier d'honneur de N.D. d'Alban fut
probablement leur propriété, heureusement cédé par la
famille Bousquet de Carménel et mis en valeur par
l'actuelle municipalité). Il y en a aussi à Montfranc et
ailleurs. Ce terme « d'Espital » accolé à un coin de
Belsert pourrait en être un lointain souvenir. La
« Chapelle de Saint Laurent » qui se trouve mentionnée
à la suite et qui fut fondée par les de Moret, possède
un pré « als abeuradou » entre le chemin de Peyreblanque
et celui de la « Fontaurie », plusieurs terres à la Valdié
et en bordure de la Voie d'Alban -- Saint Sernin.
Pareillement, elle a des immeubles en ville et hors la cité.
Qu'en reste-t-il ? ... sinon quelques écrits le rappellent.
Ainsi on peut voir sur la carte cadastrale de 1817 un
chemin des « abreurvoirs » menant à la « foutaourie\footnote{Fontaurie?} »
comme nous le disons en notre patois. Y figure encore
le « Rec de la Cappellanyé ». Une croix du même nom
dressée en 1944 en ces parages (actuellement à l'ouest
de la nouvelle gendarmerie et que fleurit joliment un
artiste en décoration florale) rappelle ce point d'histoire
en passant. Et à l'heure où notre patois (occitan) trouve
refuge à l'université, il est heureux que se